% Generated by report_gen.py (hwde v1) - do not hand-edit.
\documentclass[11pt]{article}
\newcommand\hspaper{a4paper}
\usepackage[nodiagramkit]{housestyle}
\usepackage{longtable}
\usepackage{pdfpages}
\makeatletter
\@ifundefined{KV@Gin@artifact}{\define@key{Gin}{artifact}[]{}}{}
\makeatother
\setcounter{tocdepth}{1}
\hsslug{PCB-0017-A \textperiodcentered\ stereo-class-d-amp}
\renewcommand{\sectionmark}[1]{\markboth{#1}{}}
\hssection{\leftmark}
\begin{document}
\sloppy
\hstitleblock{PCB-0017-A \textperiodcentered\ hwde v1 pipeline \textperiodcentered\ generated 2026-09-29 12:59}%
{stereo-class-d-amp \textemdash{} Highlights}%
{Phase P7; 4/5 gates pass (not passing: verify).}


\section{At a Glance}
Owner brief: rebuild and compare against a published 44 x 34 mm reference design.

\begin{itemize}
\item TPA3116D2 stereo class-D amplifier on a single 12 V supply, stereo input and output, SMD parts only,
2-layer board as small as possible. Target: match or beat the reference design's 44 x 34 mm.
\item Correctness and the TPA3116D2's thermal pad / heat-sinking come before size; say where we deviate.
\item The reference design's reported parts: TPA3116D2; 2.2 uF X7R input coupling; 100 uF + 100 nF + 1 nF supply bypass;
10 uH LC output filters per channel with damping branches; 5.6k/100k config and pull-up resistors;
wire pads J1-J4 and TP1. About 40 SMD parts, 45 footprints.
\item Source: https://jlcpcb.com/blog/gpt-6-astra-pcb-design-in-kicad
\item Stop at a DRC/ERC-clean fab package. Do NOT order. End with COMPARISON.md: our size, layers,
part count and DRC result against the reference design's.
\end{itemize}
{\small
\begin{longtable}{lp{11cm}}
\hstoprule
\hshead{Field} & \hshead{Value} \\
\hline
\endhead
part number & PCB-0017-A \\
board & stereo-class-d-amp \\
phase & P7 (Routing) \\
size & 53.5 x 36.8 mm \\
layers & 2, 2.0 oz copper (JLC2313\_1.6\_2oz) \\
parts & 39 placed parts, 16 BOM lines \\
gates & 4 of 5 recorded gates pass \\
\hline
\end{longtable}
}
\section{The Board}
\begin{center}
\includegraphics[width=0.48\textwidth]{reports/PCB-0017-A_stereo-class-d-amp_top.png}\hfill
\includegraphics[width=0.48\textwidth]{reports/PCB-0017-A_stereo-class-d-amp_bottom.png}\hfill
\\{\small top and bottom, as rendered from the board file}
\end{center}
\section{Copper Layers}
The outer copper layers as KiCad shows them, with every pad, track and zone labelled with its net (zoom in to read the small ones).
\clearpage
\subsection*{F.Cu}
\begin{center}
\includegraphics[width=\textwidth,height=0.82\textheight,keepaspectratio]{reports/layers/PCB-0017-A_stereo-class-d-amp_F_Cu.pdf}\\
{\small\texttt{reports/layers/PCB-0017-A\_stereo-class-d-amp\_F\_Cu.pdf}}
\end{center}
\clearpage
\subsection*{B.Cu}
\begin{center}
\includegraphics[width=\textwidth,height=0.82\textheight,keepaspectratio]{reports/layers/PCB-0017-A_stereo-class-d-amp_B_Cu.pdf}\\
{\small\texttt{reports/layers/PCB-0017-A\_stereo-class-d-amp\_B\_Cu.pdf}}
\end{center}
\subsection*{3D view}
\begin{center}
\includegraphics[width=0.8\textwidth]{reports/layers/PCB-0017-A_stereo-class-d-amp_iso.png}\\
{\small\texttt{reports/layers/PCB-0017-A\_stereo-class-d-amp\_iso.png}}
\end{center}
\section{What Is On It}
{\small
\begin{longtable}{p{3.6cm}p{5.4cm}p{3cm}l}
\hstoprule
\hshead{Refs} & \hshead{Part} & \hshead{Footprint} & \hshead{LCSC} \\
\hline
\endhead
C1,C2,C3,C4 & 1uF 50V X7R & C0805 & C28323 \\
C5 & 1uF 25V X7R & C0603 & C106858 \\
C6,C9,C10 & 100nF 50V X7R & C0603 & C14663 \\
C7,C8 & 1nF 50V NP0 & C0603 & C106246 \\
C11,C12 & 10uF 25V X5R & C0805 & C15850 \\
C13 & 100uF 35V polymer & CAP-SMD\_BD6.3-L6.6-W6.6-RD & C46550467 \\
C14 & 100uF 35V elec & CAP-SMD\_BD6.3-L6.6-W6.6-LS7.3-FD & C2992596 \\
C15,C16,C17,C18 & 220nF 50V X7R & C0805 & C5378 \\
C19,C20,C21,C22 & 680nF 50V X7R & C0805 & C107133 \\
C23,C24,C25,C26 & 10nF 50V X7R & C0402 & C15195 \\
D1 & SMBJ15A & SMB\_L4.6-W3.6-LS5.3-RD & C113988 \\
L1,L2,L3,L4 & 10uH 5.5A & IND-SMD\_L7.3-W6.6\_MHCI06030 & C207842 \\
R1 & 5.6k 1\% & R0603 & C23189 \\
R2,R3 & 100k 1\% & R0603 & C25803 \\
R4,R5,R6,R7 & 3.3R 1\% & R0603 & C137725 \\
U1 & TPA3118D2DAPR & HTSSOP-32\_L11.0-W6.1-P0.65-LS8.1-BL-EP & C46497 \\
\hline
\end{longtable}
}
\includepdf[pages=1]{reports/schematic.pdf}
\section{Key Decisions}
The ten that changed the board most, of 13; the full design doc has every one.
{\small
\begin{longtable}{lp{1.6cm}p{6.2cm}p{6.2cm}}
\hstoprule
\hshead{Ph} & \hshead{When} & \hshead{Decision} & \hshead{Why} \\
\hline
\endhead
P1 & 09-27 00:53 & DEVIATION: amp IC TPA3116D2DAD -\textgreater{} TPA3118D2DAP (same SLOS708 family, same 32-pin HTSSOP pinout, PowerPAD DOWN) & TPA3116D2 exists only as DAD (pad UP, TI SLOS708); its 14 C/W needs a top heatsink, which the reference design's own write-up admits is not modeled. Brief: thermal pad/heat-sinking first, SMD only. At 12 V output is rail-limited so no audio power lost; pad-... \\
P1 & 09-27 01:13 & Parts leads: TPA3118D2DAPR C46497; SLO0630H100MTT 10uH C207842 (Isat 5.5A); 680nF out filter; bulk MA25V100M6x6 polymer + RVT 100uF 35V electrolytic & P1 component-scout shortlist \\
P3 & 09-27 02:00 & C11/C12 10 uF PVCC: accept X5R CL21A106KAYNNNE (C15850) though sheets.md said X7R & caps sit below 85 C (X5R limit) at \textless{}=1.5 W design dissipation; DC-bias loss at 12 V is covered by the 100 uF bulk per pair; Basic part \\
P3 & 09-27 02:00 & R4-R7 3.3R 0603 100 mW (RC0603FR-073R3L) accepted & branch dissipation \textasciitilde{}1 mW: 10 nF is \textasciitilde{}800 ohm at 20 kHz and \textasciitilde{}40 ohm at 400 kHz against \textless{}0.3 V residual ripple \\
P2 & 09-27 02:06 & R4-R7+C23-C26 (3.3R+10nF) stay on the speaker side of the LC filter as TI's EMI C-RC snubber, not at the switch node & research: TI labels it an EMI snubber (corner \textasciitilde{}4.8 MHz), not audio damping; at the switch node 10 nF x 12 V\textasciicircum{}2 x 400 kHz = 0.58 W would overload a 0603. Filter Q 0.52-1.04 over 4-8 ohm is acceptable per SLAA701B \\
P5 & 09-27 03:27 & U1 footprint fix: the 9 EasyEDA thermal vias (0.3 drill, 0.61 pad, 1.0 pitch 3x3) had empty pad numbers (no net -\textgreater{} 9 mask-bridge + 9 clearance errors); renumbered to pad 33 (GND) and set layers {*}.Cu only (no paste/mask on the vias: bottom tented, top inside the EP mask window) & thermal vias must be GND and join the bottom pour; record class-d-amp-powerpad-thermal-via-array asks component-side capping, which is impossible inside the EP window without plugging - via-in-pad plugging is the P9 alternative; 3x3 at 1.0 mm meets the \textgreater{}=6... \\
P4 & 09-27 03:36 & review W1/W2: swap C13 to a \textgreater{}=35 V low-ESR polymer/hybrid D6.3 part and C14 to a 35 V part with a stated ESR/ripple; ADD a unidirectional TVS (SMBJ15A-class, clamp \textless{} 30 V abs max) across J1 +12V/GND & reviewer: 26 V hot-plug ring exceeds C13's 25 V and C14's ESR cannot damp the bank; TVS bounds the ring below U1's 30 V abs max. Correctness over the reference design's part count (deviation for COMPARISON) \\
P4 & 09-27 03:36 & review W3: C1-C4 2.2 uF -\textgreater{} 1 uF X7R 50 V 0805, K tolerance (TI Fig 37 value; corner 2.7 Hz at 60k) & 2.2 uF x 60k (132 ms) with no source connected reaches the 420 ms DC-detect window at tolerance corners and fault-cycles on power-up; deviation from the reference design's 2.2 uF recorded \\
P4 & 09-27 03:36 & review W4: reverse polarity stays unprotected (Q7) but J1 gets a +12V/GND silk legend; W5: C11/C12 X5R 10 uF accepted & W4 accepted risk, legend is free; W5: \textasciitilde{}4-5 uF effective at 12 V is ample beside 200 uF bulk + 100 nF/1 nF HF caps \\
P4 & 09-27 03:46 & Schematic review W1-W3 fixed: C13 35V polymer C46550467, C14 FZ 35V C2992596 (stated ESR/ripple), C1-C4 1uF X7R C28323, D1 SMBJ15A C113988 across J1 & W1 26V ring \textgreater{} 25V C13 and C14 cannot damp it -\textgreater{} TVS clamps; W2 RVT had no ESR spec; W3 132ms input tau vs 420ms DC-detect -\textgreater{} 60ms \\
\hline
\end{longtable}
}
\section{How the Run Went}
37 recorded decisions from P0 to P7, 2026-09-27T00:36 to 2026-09-27T05:25. The run went back to an earlier phase 3 times.
\begin{itemize}
\item P3 back to P2 (02:06): R4-R7+C23-C26 (3.3R+10nF) stay on the speaker side of the LC filter as TI's EMI C-RC snubber, not at the switch node
\item P5 back to P4 (03:36): review W1/W2: swap C13 to a \textgreater{}=35 V low-ESR polymer/hybrid D6.3 part and C14 to a 35 V part with a stated ESR/ripple; ADD a unidirectional TVS (SMBJ15A-class, clamp \textless{} 30 V abs max) across J1 +12V/GND
\item P7 back to P6 (05:25): spread L3/L4 up and L1/L2 down \textasciitilde{}2.5 mm (grow outline height if needed) to fit the C15-C18 bootstrap caps in two columns and route OUT nets on F.Cu
\end{itemize}
\section{Checks and Cost}
{\small
\begin{longtable}{llrr}
\hstoprule
\hshead{Gate} & \hshead{Status} & \hshead{Attempts} & \hshead{Failing} \\
\hline
\endhead
erc & pass & 2 & 0 \\
place & pass & 31 & 0 \\
drc\_routed & pass & 14 & 0 \\
verify & fail & 5 & 8 \\
dfm & pass & 8 & 0 \\
\hline
\end{longtable}
}
Estimated JLCPCB cost for 5 assembled boards: \$61.09 (\$12.22 each). It is an estimate from the price table, not a quote.

Model cost of generating this board: \$97.57.

The full design doc (stereo-class-d-amp-design-full) has every decision, the renders and the run's history.

\end{document}
